% Section 6.6 — Classification Error Analysis
% Error analysis and confusion patterns.

\section{Classification Error Analysis}
\label{sec:classification-errors}

This section analyses the classification errors made by both the proposed method and the CE baseline on the SDI-4 test partition.

\subsection{Confusion Pattern}
\label{sec:confusion-pattern}

The primary sources of confusion are expected to occur between classes that share visual characteristics. Specific patterns include:

\begin{itemize}
    \item \textbf{WSSV vs.~WSSV\_BG}: Images exhibiting co-infection may be mistaken for single-infection WSSV when the black gill symptoms are subtle or absent.
    \item \textbf{Healthy vs.~disease classes}: False positives on Healthy images (predicting a disease class for a healthy specimen) and false negatives (missing a diseased specimen) carry different operational implications.
    \item \textbf{BG vs.~WSSV\_BG}: The presence of black gill symptoms in both classes creates a natural boundary ambiguity.
\end{itemize}

Detailed confusion matrices are available from the retained experiment logs and will be reported in the final thesis.

\subsection{Error Reduction}
\label{sec:error-reduction}

The overall accuracy improvement from 89.02\% (CE baseline) to 91.33\% (proposed) corresponds to approximately 4 fewer misclassified images out of the 173-image test set. The component-level contribution to error reduction will be analysed when the individual component logs are extracted.